\section*{Hoofdstuk \ref{ch:b2:alpine-orogen} --- Anatomie van een gebergtegordel: de Alpen}

\begin{solution}{exo:b2:alpine-orogen:1}
(a) $e = 1 - 27.0/36.0 = 0.25$, dat wil zeggen \qty{25}{\percent}; $S = 36.0 - 27.0 =
\qty{9.0}{km}$. (b) De laag moet haar lengte hebben behouden (geen uitrekking of verdunning langs de laag),
de doorsnede moet evenwijdig aan de transportrichting liggen, en de pinlijnen moeten plaatsen markeren
waar de lagen niet langs elkaar zijn geschoven.
\end{solution}

\begin{solution}{exo:b2:alpine-orogen:2}
(a) Een \omterm{def:b2:alpine-orogen:klippe}{klippe}: de gneis is een geïsoleerd restant van een \omterm{def:b2:alpine-orogen:nappes}{dekblad} dat op zijn \omterm{def:b2:alpine-orogen:nappes}{autochtoon} rust;
erosie heeft de rest van de plaat rondom verwijderd. (b) Een \omterm{def:b2:alpine-orogen:klippe}{tektonisch venster}: erosie heeft
het gneisdekblad doorsneden en de jongere gesteenten onder zijn basale overschuiving blootgelegd, die
daarom aan alle kanten van het venster weg helt.
\end{solution}

\begin{solution}{exo:b2:alpine-orogen:3}
(a) Ultramafisch complex (de bovenmantel); gabbrocomplex (de onderste oceanische korst);
complex van parallelle \omterm{def:b2:magma-emplacement:bodies}{gangen} (de middelste korst); vulkanisch complex, met kussens (de bovenste korst).
(b) Het complex van parallelle \omterm{def:b2:magma-emplacement:bodies}{gangen}. (c) De meeste korst van de wijde oceanen wordt door subductie verbruikt;
de \omterm{def:b2:alpine-orogen:ophiolite}{ofiolieten} die op het land overleven, ontstonden vooral in kleine oceaanbekkens of tussen
eilandbogen, die van de Tethys waar een continent dunner werd tot er een kleine oceaan openging.
\end{solution}

\begin{solution}{exo:b2:alpine-orogen:4}
(a) \omterm{def:b2:alpine-orogen:nappes}{Dekbladen} van de Monte Rosa en de Grote Sint-Bernhard: de Europese kant (het Briançonnais en de
Europese passieve rand); Zermatt--Saas en de Combinzone: de Piemont--Ligurische oceaan; systeem van de Dent
Blanche: de rand van Apulië. (b) Apulië, dat bovenop eindigt. (c) De bovenplaat wordt
niet gesubduceerd, en de schijven die van de onderplaat worden afgeschraapt, stapelen zich eronder op in de volgorde waarin de
onderplaat aankwam: verre oceaan eerst, nabije rand als laatste en het laagst. Latere plooiing van de stapel,
of overschuivingen buiten de normale volgorde, kunnen de volgorde verstoren.
\end{solution}

\begin{solution}{exo:b2:alpine-orogen:5}
(a) $2\pi G\,\Delta\rho = 2\pi\times6.674\times10^{-11}\times500 =
\qty{2.097e-7}{s^{-2}}$; $h = 1.70\times10^{-3}/2.097\times10^{-7} = \qty{8.1}{km}$.
(b) $2.097\times10^{-7}\times25\,000 = \qty{5.24e-3}{m.s^{-2}}$, ongeveer
\qty{-524}{mGal}. (c) De wortel is geen oneindige plaat: het is een lichaam van enkele tientallen kilometers
diep en slechts ongeveer zo breed als de gordel, zodat het minder trekt aan een punt aan het oppervlak dan een
plaat van dezelfde dikte die tot in het oneindige doorloopt; dichte lichamen zoals de Ivreazone
voegen positieve anomalieën toe, en het dichtheidscontrast op diepte kan kleiner zijn dan aangenomen.
\end{solution}

\begin{solution}{exo:b2:alpine-orogen:6}
(a) $164 - 44 = \qty{120}{Myr}$. (b) $z = P/(\bar\rho g)$:
$2.8\times10^9/(3000\times9.8) = \qty{95.2}{km}$; $2.8\times10^9/(3300\times9.8) =
\qty{86.6}{km}$. (c) ``Meer dan \qty{90}{km}'' ligt binnen dit bereik: de diepte hangt evenzeer af van
de gemiddelde dichtheid die voor de bovenliggende gesteenten wordt gekozen als van de druk.
\end{solution}

\begin{solution}{exo:b2:alpine-orogen:7}
(a) Een flank met lengte $L$ die onder $\theta$ helt, overspant horizontaal $L\cos\theta$; alle flanken
hebben dezelfde helling, dus $l = l_0\cos\theta$ en $e = 1 - l/l_0 = 1 - \cos\theta$.
(b) $\cos\theta = 0.70$: $\theta = \qty{45.6}{\degree}$. (c) $1 - \cos\qty{40}{\degree} =
\qty{23.4}{\percent}$, zoals in de chevronfiguur; $1 - \cos\qty{60}{\degree}
= \qty{50}{\percent}$.
\end{solution}

\begin{solution}{exo:b2:alpine-orogen:8}
(a) $l_0 = 118/2.0 = \qty{59.0}{km}$; $S = 59.0 - 42 = \qty{17.0}{km}$; $e = 17.0/59.0 =
\qty{28.8}{\percent}$. (b) $A_e = 118 - 42\times2.0 = \qty{34.0}{km^2}$; $d = 34.0/17.0 =
\qty{2.0}{km}$. (c) Ja: het \omterm{def:b2:alpine-orogen:decollement}{décollement} ligt aan de basis van het dek van \qty{2.0}{km}, onder
zijn top; het dek is langs zijn basis losgemaakt.
\end{solution}

\begin{solution}{exo:b2:alpine-orogen:9}
(a) De Airy-wortel is $4.5h$: op \qty{46.5}{\degree} NB, $4.5\times2.74 = \qty{12.3}{km}$, een
Moho op \qty{42.3}{km}; op \qty{45.55}{\degree} NB, $4.5\times0.65 = \qty{2.9}{km}$, een Moho
op \qty{32.9}{km}. (b) Onder het hoge terrein ligt de Moho \qty{6.0}{km} ondieper dan Airy
voorspelt; onder de rand van de Povlakte ligt hij \qty{28.3}{km} dieper. (c) De wortel is niet
lokaal gecompenseerd: hij is naar het zuiden verschoven en groter dan het reliëf vereist.
De plaat buigt als geheel onder de lasten die erop rusten (de regionale ondersteuning uit een hoofdstuk van
Boek~5), en een van die lasten, het Ivrealichaam, is helemaal geen topografie.
\end{solution}

\begin{solution}{exo:b2:alpine-orogen:10}
(a) $5.1/11.3 = 0.45$: minder dan de helft. (b) $5.1 + 6.2 = \qty{11.3}{km}$. De top van het
Jura ligt boven het hogere onthechtingsvlak; de \qty{6.2}{km} die langs dat onthechtingsvlak werden meegevoerd,
gingen onder de plooien van de doorsnede door en werden verbruikt in structuren verder naar het noorden, zodat de
bovenste laag binnen de doorsnede nooit met dat bedrag verkortte. (c) Eén gidslaag registreert
alleen de verkorting van de lagen waartoe ze behoort; elk onthechtingsvlak erboven of eronder kan
verplaatsing de doorsnede in of uit voeren. Balanceer meerdere lagen, en controleer met oppervlakten.
\end{solution}

\begin{solution}{exo:b2:alpine-orogen:11}
(a) Voor een verplaatsing van \qty{100}{km}, $\sqrt{200^2 + 100^2} = \qty{223.6}{km}$; voor
\qty{50}{km}, $\sqrt{200^2 + 50^2} = \qty{206.2}{km}$; voor \qty{150}{km},
$\sqrt{200^2 + 150^2} = \qty{250.0}{km}$. (b) $\arctan(100/200) = \qty{26.6}{\degree}$ ten opzichte van
noord--zuid, een convergentielijn dicht bij de N\qty{150}{\degree}E die Rosenberg en
collega's opgeven. (c) De verplaatsing is onzeker met $\pm\qty{50}{km}$, wat de totale
convergentie met ongeveer \qty{20}{km} in beide richtingen verschuift; het derde cijfer van het resultaat ligt ver onder
die onzekerheid.
\end{solution}

\begin{solution}{exo:b2:alpine-orogen:12}
(a) $210/12.9 = \qty{16.3}{Myr}$ tot $210/8.5 = \qty{24.7}{Myr}$ (\unit{km} gedeeld door
\unit{mm.yr^{-1}} geeft miljoenen jaren). (b) Meer dan $4\times60 = \qty{240}{km}$; minstens
$240 - 60 = \qty{180}{km}$ van het bekken ligt nu onder de Alpiene overschuivingen of is
geërodeerd. (c) Het bekken ontstaat vóór de overschuivingswig en migreert ermee: zijn binnenste
deel wordt steeds overreden door de overschuivingen die het belasten, en zijn oudste vulling raakt gevangen
in de overschuivingen, zoals de Subalpiene \omterm{def:b2:alpine-orogen:foreland}{Molasse}.
\end{solution}

\begin{solution}{pb:b2:alpine-orogen:1}
\noindent\textbf{Deel I.}

\textbf{1.} $3.0/\cos\qty{35}{\degree} = \qty{3.66}{km}$; $2.5/\cos\qty{50}{\degree} =
\qty{3.89}{km}$; $4.0/\cos\qty{40}{\degree} = \qty{5.22}{km}$.

\textbf{2.} $8.00 + 2\times3.66 + 5.00 + 2\times3.89 + 6.0/\cos\qty{10}{\degree} +
2\times5.22 + 10.0/\cos\qty{5}{\degree} = 8.00 + 7.32 + 5.00 + 7.78 + 6.09 + 10.44 + 10.04 =
\qty{54.67}{km}$; met het herhaalde stuk is $l_0 = \qty{60.7}{km}$.

\textbf{3.} De laag onder de overschuivingsplaat is dezelfde laag, door de overschuiving onder de plaat
gebracht: vóór de overschuiving lag ze in het verlengde van de rest. Als men haar weglaat, wordt de
laag te kort hersteld.

\textbf{4.} $e = 1 - 48.0/60.7 = 0.209$, dat wil zeggen \qty{20.9}{\percent}; $S = 60.7 - 48.0 =
\qty{12.7}{km}$.

\textbf{5.} $60.7 - 60.2 = \qty{0.5}{km}$, $0.5/60.7 = \qty{0.8}{\percent}$: binnen de
precisie van de meting; de doorsnede is aanvaardbaar.

\textbf{6.} $60.7 - 54.0 = \qty{6.7}{km}$, \qty{11}{\percent}: de twee lagen van één plaat
kunnen niet tot zulke verschillende lengten herstellen. De doorsnede zou zoals getekend niet toelaatbaar zijn: een
structuur die ontbreekt of verkeerd is getekend, een pinlijn op de verkeerde plaats, beweging uit het vlak van
de doorsnede, of een onthechtingsvlak tussen de twee lagen.

\noindent\textbf{Deel II.}

\textbf{7.} $146/2.4 = \qty{60.8}{km}$.

\textbf{8.} Het verschilt \qty{0.1}{km} van \qty{60.7}{km}: lengte en oppervlakte stemmen overeen.

\textbf{9.} $A_e = 146 - 48.0\times2.4 = 146 - 115.2 = \qty{30.8}{km^2}$.

\textbf{10.} $d = 30.8/12.7 = \qty{2.43}{km}$, gelijk aan de \qty{2.4}{km} van het dek binnen
de precisie van de gegevens: het \omterm{def:b2:alpine-orogen:decollement}{décollement} is de basis van het dek.

\textbf{11.} $160/2.4 = \qty{66.7}{km}$, \qty{6.0}{km} meer dan de lengte langs de laag. Verdikking
van de evaporieten aan de basis van het dek, die in de kernen van anticlinalen vloeien, voegt oppervlakte toe
zonder lengte toe te voegen; oppervlaktebalans met een constante dikte faalt dan.

\textbf{12.} Het dek bewoog over een \omterm{def:b2:alpine-orogen:decollement}{décollement} aan zijn basis en de sokkel was er niet
bij betrokken. Onder de gordel zou de top van de sokkel vlak moeten lopen of zacht en
gelijkmatig moeten hellen, ongebroken door de overschuivingen erboven.

\noindent\textbf{Deel III.}

\textbf{13.} $22 + 18 + 25 = \qty{65}{km}$; $S = 65 - 15 = \qty{50}{km}$.

\textbf{14.} $e = 50/65 = \qty{76.9}{\percent}$, tegen \qty{20.9}{\percent} in de
voorlandgordel: de \omterm{def:b2:alpine-orogen:nappes}{dekbladen} zijn veel sterker verkort.

\textbf{15.} $30 - 22 = \qty{8}{km}$ ($e = \qty{26.7}{\percent}$). De sokkel zelf is
verkort: \omterm{def:b2:alpine-orogen:skin}{dikhuidige tektoniek}.

\textbf{16.} Voorlandgordel $12.7/(11 - 4) = \qty{1.81}{km/Myr}$, dat wil zeggen
\qty{1.81}{mm} per jaar; \omterm{def:b2:alpine-orogen:nappes}{dekbladen} $50/(30 - 20) = \qty{5.0}{mm}$ per jaar.

\textbf{17.} De \omterm{def:b2:alpine-orogen:nappes}{dekbladen} (30 tot \qty{20}{Ma}) werden vervormd vóór de voorlandgordel (11 tot
\qty{4}{Ma}): de vervorming migreerde naar buiten, naar het voorland. Dat is de volgorde van de
Alpen: de Glarus-overschuiving in het Oligoceen en Mioceen, de externe massieven in het Mioceen,
de Jura in het Laat-Mioceen en Plioceen.

\textbf{18.} $12.7 + 50 + 8 = \qty{70.7}{km}$.

\noindent\textbf{Deel IV.}

\textbf{19.} $\qty{0.5}{cm.yr^{-1}}\times\qty{40}{Myr} = 0.005\times4\times10^7 =
\qty{2.0e5}{m} = \qty{200}{km}$.

\textbf{20.} $200 + 18 = \qty{218}{km}$.

\textbf{21.} Alleen wat bewaard is of wordt afgeleid, telt mee: delen van de oceaan kunnen
vóór \qty{85}{Ma} of sneller dan de gemiddelde snelheid zijn gesubduceerd, ofiolietschijven kunnen
geërodeerd of uit de doorsnede weggevoerd zijn, en convergentie die schuin op de doorsnede staat, voert
oceaanbodem aan die de doorsnede niet ziet.

\textbf{22.} Interne zone $60 - 20 = \qty{40}{km}$ ($e = \qty{66.7}{\percent}$); retrowig
$70 - 48 = \qty{22}{km}$ ($e = \qty{31.4}{\percent}$).

\textbf{23.} $12.7 + 50 + 8 + 40 + 22 = \qty{132.7}{km}$; binnen het bereik van de schattingen
voor de Centrale Alpen, ongeveer 125 tot \qty{160}{km} botsingsverkorting (voordat de
horizontale verschuiving wordt opgeteld).

\textbf{24.} Herstelde breedte $153 + 132.7 = \qty{285.7}{km}$; herstelde korst
$285.7\times30 = \qty{8571}{km^2}$; huidige korst $153\times45 = \qty{6885}{km^2}$. Ontbrekend:
$8571 - 6885 = \qty{1686}{km^2}$, \qty{19.7}{\percent} van de herstelde korst; ze is
geërodeerd (en in de \omterm{def:b2:alpine-orogen:foreland}{voorlandbekkens} afgezet) of met de onderkorst in de
mantel gesubduceerd. Zonder die verliezen zou de korst gemiddeld $8571/153 = \qty{56}{km}$ dik zijn.
De totale botsingsverkorting over de doorsnede is
\textbf{\qty[text-series-to-math=true]{133}{km}}.
\end{solution}
